\section{Threats to Validity}
\label{sec:threats}

\paragraph{Single GPU.}
All measurements are on one \gpuModel{}. We have no H100 comparison, so the
absolute latencies and the crossover constants carry an unknown
hardware-scaling term. The cost \emph{laws} (linear vs.\ staircase) should
transfer---they are structural---but the fitted constants ($\arBase$,
$\arSlope$, $\dgBase$, $\dgStep$) are A100-specific.

\paragraph{Single model family for the staircase.}
The staircase cost law is established on DiffusionGemma-26B-A4B and confirmed
across two corpora (BGL: step $= 118.5$\,ms; HDFS: step $= 120.5$\,ms). The
LLaDA-8B arm under dInfer was measured but is \emph{inconclusive}: dInfer's
prefix/vicinity KV-cache errors on prompts $\geq 632$ tokens force a
cache-less mode that inflates absolute latency by $\sim 30\times$ and
introduces variance that obscures any block-length-based stepping. The
\emph{form} of the law should hold for any canvas-bound diffusion model, but
LLaDA's confirmation requires a production-grade runtime with working KV
cache, which does not exist at time of writing.

\paragraph{Long-context memory.}
We attempted to measure AR-1.5B's GPU memory at context lengths up to 32k
with \texttt{enforce\_eager}, but \texttt{torch.cuda.max\_memory\_allocated}
returns 0 because vLLM manages its own memory pool internally. The latency
data confirms AR's input-length sensitivity (43\,ms at 1k context $\to$
436\,ms at 16k), but the memory footprint question remains open. DG's
absence of a growing KV cache may give it a flatter memory curve at long
context, but this was not measured.

\paragraph{Two-variable scope.}
The cost law is two-variable: $T_{DG}(L_{out}, L_{in})$. The per-event
crossover $L^\ast=\crossoverL$~tokens is derived from the one-variable
projection at the corpus's few-shot input length ($\sim\!1047$ tokens). At
longer inputs the diffusion model's steeper input slope (\dgInputSlope{} vs
\arInputSlope{}~ms/prompt-token) raises the base and pushes $L^\ast$ upward.
The one-variable crossover is valid only near the input length at which it was
calibrated and must not be extrapolated to arbitrary prompt lengths.

\paragraph{Not size-matched.}
The compute-matched pair (AR-3B vs DG-26B-A4B; 3B vs 4B active parameters) is
matched on active compute, not on size. The quality-matched AR model
(Qwen2.5-1.5B, $\sim$3~GB) is $17\times$ smaller than DiffusionGemma (52~GB).
The per-event latency advantage the diffusion model holds at batch-1 does not
survive this footprint comparison: a $17\times$ smaller model that matches
quality is the deployment winner whenever memory or cost matters, which is
whenever it is deployed.

\paragraph{$K$ is pinned.}
$K=8$ denoising steps is pinned throughout. The base is $K$-dependent: at
$K=32$ it was \dgOldBase{}~ms. A different $K$ changes both the base and the
crossover ($L^\ast=\crossoverL$ at $K=8$ vs $65$ at $K=32$). The decision
procedure takes $K$ as an input so the crossover recomputes, but every number
in this paper assumes $K=8$ unless stated.
